The cryptographic systems discussed in this report rely on a small collection of mathematical ideas from number theory and computational hardness. This section presents only the background necessary to understand ElGamal and RSA. It introduces the concepts that directly support the security arguments and correctness proofs used later.

\subsection{Modular Arithmetic}

Modular arithmetic studies integers with respect to a fixed modulus $n \geq 2$. For integers $a$ and $b$, one writes
\[
a \equiv b \pmod{n}
\]
when $n$ divides $a-b$. In other words, $a$ and $b$ leave the same remainder when divided by $n$. Computation modulo $n$ allows large values to be reduced into a finite set of residue classes
\[
\{0,1,2,\dots,n-1\}.
\]
This structure is fundamental in cryptography because encryption and decryption algorithms often involve repeated multiplication and exponentiation modulo a large integer. \textbf{Modular arithmetic} is therefore the basic language in which both ElGamal and RSA are expressed.

For example, if $17 \equiv 5 \pmod{12}$, then both integers represent the same residue class modulo $12$. Likewise, arithmetic operations are preserved under congruence. If
\[
a \equiv b \pmod{n}
\quad \text{and} \quad
c \equiv d \pmod{n},
\]
then
\[
a+c \equiv b+d \pmod{n}
\]
and
\[
ac \equiv bd \pmod{n}.
\]
This property makes modular arithmetic especially suitable for algorithmic computation in cryptographic constructions.

\subsection{Euler's Totient Function}

Euler's totient function, denoted by $\varphi(N)$, counts the number of integers in the set $\{1,2,\dots,N-1\}$ that are relatively prime to $N$. This function plays a central role in RSA. In the important special case where
\[
N = pq
\]
for two distinct primes $p$ and $q$, the totient value is
\[
\varphi(N) = (p-1)(q-1).
\]
The reason is that among the integers from $1$ to $N-1$, the only values not relatively prime to $N$ are those divisible by $p$ or $q$.

This formula is essential because \textbf{RSA key generation} depends on choosing a public exponent $e$ and a private exponent $d$ such that
\[
ed \equiv 1 \pmod{\varphi(N)}.
\]
The existence of this modular inverse requires
\[
\gcd(e,\varphi(N)) = 1.
\]

\subsection{Euler's Theorem}

Euler's theorem is a basic result connecting modular exponentiation and the totient function. It states that if
\[
\gcd(x,N) = 1,
\]
then
\[
x^{\varphi(N)} \equiv 1 \pmod{N}.
\]
This theorem generalizes Fermat's little theorem and provides the \textbf{core mathematical justification for RSA decryption}. If the exponents $e$ and $d$ satisfy
\[
ed = k\varphi(N) + 1
\]
for some integer $k$, then
\[
x^{ed} = x^{k\varphi(N)+1} = \left(x^{\varphi(N)}\right)^k x \equiv x \pmod{N},
\]
assuming $\gcd(x,N)=1$. This is the central idea behind the correctness of the RSA trapdoor permutation.

\subsection{Discrete Logarithm Problem}

The ElGamal cryptosystem is based on the difficulty of solving the \textbf{discrete logarithm problem} in an appropriate cyclic group. Let $G$ be a cyclic group generated by an element $g$. Given an element
\[
h = g^a,
\]
the discrete logarithm problem asks for the recovery of the exponent $a$. Although exponentiation is efficient, the inverse operation is believed to be computationally hard in carefully chosen groups of large order.

This asymmetry between easy forward computation and hard inversion is the same general pattern that appears throughout public key cryptography. In ElGamal, the public key contains a group element of the form $g^a$, while the private key is the exponent $a$ itself. Security depends on the assumption that an adversary cannot efficiently recover $a$ from the public information alone.

\subsection{Computational Diffie--Hellman Assumption}

The \textbf{Computational Diffie--Hellman (CDH) assumption} formalizes the difficulty of combining exponentiated group elements into a shared secret. Given
\[
g, \quad g^a, \quad g^b,
\]
the task is to compute
\[
g^{ab}.
\]
The CDH assumption states that no efficient algorithm can perform this computation with non-negligible probability in the groups used by the scheme.

This problem is closely related to the discrete logarithm problem, but it is not identical to it. Even if an adversary cannot recover $a$ or $b$ individually, the question remains whether the shared value $g^{ab}$ can be computed directly. In many cryptographic settings, CDH is treated as a natural hardness assumption for key agreement and related constructions.

\subsection{Decisional Diffie--Hellman Assumption}

The \textbf{Decisional Diffie--Hellman (DDH) assumption} is stronger than CDH and is particularly important for the security of ElGamal encryption. Given a tuple
\[
(g, g^a, g^b, T),
\]
the challenge is to decide whether
\[
T = g^{ab}
\]
or whether $T$ is simply a random group element independent of $a$ and $b$. Under the DDH assumption, no efficient adversary can distinguish these two cases with significant advantage.

This decisional formulation matters because semantic security requires that ciphertext components appear random to an attacker. In ElGamal, the value $h^r = g^{ar}$ is hidden inside the ciphertext. If the tuple formed from the public values and this hidden term is computationally indistinguishable from random, then the encrypted message is also protected against efficient adversaries.

\subsection{Factoring Assumption}

RSA is built on a different hardness assumption. Instead of relying on discrete logarithms in groups, it relies on the presumed difficulty of factoring a large composite integer
\[
N = pq,
\]
where $p$ and $q$ are large primes. The factoring assumption states that, given only $N$, no efficient algorithm can recover its prime factors when the modulus is chosen sufficiently large and properly generated.

This \textbf{factoring assumption} is important because knowledge of $p$ and $q$ immediately reveals
\[
\varphi(N) = (p-1)(q-1),
\]
which in turn allows the computation of the private exponent $d$ from the public exponent $e$. Thus, the security of RSA is connected to the difficulty of reconstructing the hidden factorization of $N$.

Taken together, these mathematical concepts establish the background for the two major systems studied in this chapter. \textbf{ElGamal} depends on hardness assumptions related to exponentiation in cyclic groups, especially DDH, whereas \textbf{RSA} depends on arithmetic modulo a composite integer and the difficulty of factoring that modulus.
